% day2_seg4_numerical/roundoff.tex
\begin{frame}
  \onslide<1->{
	\frametitle{Round-Off Errors} 
	\framesubtitle{Joseph Form}
}
\begin{itemize}
  \onslide<2->{
 \item Computer round-off in floating point arithmetic
}
  \onslide<3->{
 \item Unit round-off value is largest number $\eps$ such that $1+\eps=1$
}
  \onslide<4->{
 \item In Matlab, the global symbol "eps" stores the unit round-off value
}
  \onslide<5->{
  \item Round-off errors generally plague the classical covariance update algorithm
\[\covupdateeq\]
\begin{itemize}
  \item Alternative, called the Joseph form, is more robust:
\[\josephcovupdateeq\]
\end{itemize} 
}
\end{itemize}
\end{frame}

\begin{frame}
  \onslide<1->{
	\frametitle{Round-Off Errors} 
	\framesubtitle{De Vries Mechanization}
}
\begin{itemize}
  \onslide<2->{
\item   This form in turn leads to the De Vries mechanization of the
Kalman gain equation.  To see this, first expand
the Joseph form (omitting subscripts for clarity):

}
  \onslide<3->{
\begin{align*}
   P(+)&= P(-)-KHP(-)-P(-)H\trans K\trans+ KHP(-)H\trans K\trans +
   KRK\trans.
\end{align*}
}
\end{itemize}
\end{frame}

\begin{frame}
  \onslide<1->{
	\frametitle{Round-Off} 
	\framesubtitle{De Vries}
}
\begin{itemize}
  \onslide<2->{
\item   Define
\begin{equation}\label{Feq}
  F=PH\trans
\end{equation}
and
\begin{equation}\label{Deq}
  D=HF+R.
\end{equation}

}
  \onslide<3->{
\item   We note that $D$ is symmetric, because
\begin{align*}
  D\trans &= [HF+R]\trans \\
  &= [HPH\trans + R]\trans \\
  &= [HPH\trans] \trans + R\trans \\
  &= [H\trans]\trans P\trans H\trans + R\trans \\
  &= HPH\trans + R & \text{(because $P$ and $R$ are symmetric)}\\
  &= HF + R \\
  &= D.
\end{align*}
}
\end{itemize}
\end{frame}

\begin{frame}
  \onslide<1->{
	\frametitle{Round-Off} 
	\framesubtitle{De Vries}
}
\begin{itemize}
  \onslide<2->{
\item   The gain in standard Kalman gain equation
\[\kalgaineq\] can be written as
\[K=FD^{-1}.\]
}
  \onslide<3->{
\item   Using Eqs.~(\ref{Feq}) and ~(\ref{Deq}), we compute
\begin{align*}
  P(+) &= P(-) - KF\trans - FK\trans + KDK\trans \\
  &= P(-)-K\left[F\trans - \frac{DK\trans}{2} \right] - \left[
  F-\frac{KD}{2}\right]K\trans.
\end{align*}

}
  \onslide<4->{
\item   Letting $G=F-\dfrac{KD}{2}$, we obtain
\begin{equation}\label{dvcovupdateeq}
  \dvcovupdateeq.
\end{equation}

}
\end{itemize}
\end{frame}

\begin{frame}
  \onslide<1->{
	\frametitle{Round-Off} 
	\framesubtitle{De Vries}
}
\begin{itemize}
  \onslide<2->{
\item   The De Vries mechanization implements Eq.~(\ref{dvcovupdateeq}) by
subtracting the $ij^{th}$ element of $GK\trans$ from both the
$ij^{th}$ and $ji^{th}$ elements of $P(-)$. 
}
  \onslide<3->{
\item   The De Vries form of
the covariance update equation is valid for any arbitrary gain
matrix $K$, not just the optimal gains computed from
Eqs.~(\ref{Feq}) and (\ref{Deq}).

}
\end{itemize}
\end{frame}
